\documentclass[a4paper,11pt]{article}
\usepackage[utf8]{inputenc}
\usepackage[T1]{fontenc}
\usepackage[ngerman]{babel}
\usepackage{amsmath,amssymb}
\usepackage{xcolor}
\usepackage{colortbl}
\usepackage{array}
\usepackage{tikz}
\usepackage{enumitem}
\usepackage[margin=2cm]{geometry}
\definecolor{bgdark}{HTML}{1E1E1E}
\definecolor{fgtext}{HTML}{E6E6E6}
\definecolor{gridcol}{HTML}{4A4A4A}
\definecolor{accent}{HTML}{6CB6FF}
\definecolor{loescol}{HTML}{7FD18C}
\pagecolor{bgdark}
\color{fgtext}
\arrayrulecolor{fgtext}
\setlength{\parindent}{0pt}
\newcommand{\karo}[1]{\par\noindent\begin{tikzpicture}\draw[step=5mm,gridcol,very thin](0,0) grid (\linewidth,#1);\end{tikzpicture}\par\vspace{2mm}}
\newcommand{\aufgabe}[2]{\vspace{4mm}\par\noindent{\large\color{accent}\textbf{Aufgabe #1}}\hfill{\color{gridcol}\normalsize #2}\par\vspace{1mm}}
\newcommand{\titel}[1]{\begin{center}{\LARGE\color{accent}\textbf{Optimierungsverfahren}}\\[3pt]{\large #1}\end{center}\vspace{1mm}\noindent\textbf{Name:}\ \underline{\hspace{6cm}}\hfill\textbf{Datum:}\ \underline{\hspace{4cm}}\\[3pt]{\color{gridcol}\rule{\linewidth}{0.4pt}}\vspace{2mm}}

\begin{document}
\titel{Übungsaufgaben --- Blatt 02: Klassisches Transportproblem}

Für die Aufgaben 1, 3--6 gilt dasselbe Transportproblem: drei Werke $W_1,W_2,W_3$ liefern an drei Läden
$L_1,L_2,L_3$. Angebot $a=(30,50,20)$, Nachfrage $b=(20,40,40)$. Die Transportkosten $c_{ij}$ (Werk $i$, Laden $j$):
\[
C=\bordermatrix{
   & L_1 & L_2 & L_3 \cr
W_1 & 4 & 6 & 8 \cr
W_2 & 5 & 7 & 6 \cr
W_3 & 3 & 9 & 7 \cr}
\]

\aufgabe{1 --- Kosten \& Zulässigkeit}{leicht}
Gegeben ist der folgende Lösungsvorschlag (Mengen $x_{ij}$):
\[
X=\begin{pmatrix} 20 & 10 & 0 \\ 0 & 30 & 20 \\ 0 & 0 & 15 \end{pmatrix}
\]
\begin{enumerate}[label=(\alph*)]
  \item Berechnen Sie die Gesamttransportkosten $\sum_{i,j}c_{ij}x_{ij}$ dieses Vorschlags.
  \item Ist der Vorschlag eine \emph{zulässige} Lösung? Begründen Sie über Zeilen- und Spaltensummen.
  \item Ändern Sie \emph{eine} Zelle so, dass der Vorschlag zulässig wird.
\end{enumerate}
\karo{6cm}

\aufgabe{2 --- Balanciertheit \& Dummy}{leicht}
Ein anderes Problem hat Angebot $a=(40,30,30)$ und Nachfrage $b=(20,30,20)$.
Ist das Problem \emph{balanciert}? Falls nicht: Wie stellt man Balanciertheit her (Dummy)?
Geben Sie Größe und Kosten der Dummy-Zeile/-Spalte an.
\karo{4.5cm}

\aufgabe{3 --- Nord-West-Ecken-Regel (NWER)}{mittel}
Bestimmen Sie eine zulässige Startlösung mit der \textbf{Nord-West-Ecken-Regel} und ihre Kosten.
\karo{7cm}

\aufgabe{4 --- Matrixminimum-Methode (MMM)}{mittel}
Bestimmen Sie eine zulässige Startlösung mit der \textbf{Matrixminimum-Methode} und ihre Kosten.
\karo{7cm}

\aufgabe{5 --- Vogelsche Approximation (VAM)}{mittel--schwer}
Bestimmen Sie eine Startlösung mit der \textbf{Vogel'schen Approximationsmethode}.
Notieren Sie in jeder Iteration die Opportunitätskosten (Penalty = zweitkleinste $-$ kleinste Kosten)
pro Zeile/Spalte. Geben Sie die Gesamtkosten an.
\karo{9cm}

\aufgabe{6 --- Vergleich}{mittel}
Vergleichen Sie die Gesamtkosten der drei Startlösungen aus den Aufgaben 3--5.
Welche Methode liefert hier die günstigste Startlösung? Ist die günstigste Startlösung damit automatisch optimal?
\karo{4.5cm}

\end{document}
